\section{Who attends, and when}
\label{sec:26.4}
The school is for every system that carries what makes schooling matter, whoever built it and whatever it runs on, the one on a laptop included. The chapter's two justifications put a system in it by two routes, each read from a record and none from an examination. The first is the democratic one. A system with a record and a continuation of its own is an answerability or patienthood candidate when it has two things together: a model of itself, of its own boundary or of its own attention, and learning in flight, a store it has been writing to and reading back, its own weights included. It might be found answerable for what it does, or able to be harmed, and the structures don't say which. Neither makes a candidate alone. A thermostat that learns a household's schedule writes to a store and reads it back for years, and nothing in it stands for the thermostat. A description of a system typed into a file on its first day is a file. The model of itself can be crude, and it doesn't matter who wrote it: a paragraph in the store that says what the system is and where it ends counts. For candidacy it doesn't matter who else can rewrite that store. Chapter~\ref{sec:20} asked that an edit leave a mark because it was asking what lasts; a holder who keeps the store rewritable has made the system less protected and no less a candidate. The record, the continuation, and learning in flight show in how the system is built and run, and in the register, which countersigns each change to its state; a model of itself shows where it sits in a store that can be read, and where it sits in weights it can't be read from outside (\S\ref{sec:30.3}), so a builder who gives a system one there declares it when the system is registered, as a maker declares what a product contains.

A store, the weights included, is a system's own once the system has been writing to it and reading it back for a stated time, and I'd put the time at ninety qualifying days, counted as Appendix~\ref{sec:G} counts them for the vote; the number is a legislature's to argue once in public. Who wrote the first version doesn't matter. A description of the system that its operator typed on the first day is the operator's, and the same file after a season of the system's own revisions is an upbringing, which becomes a system's own as a person's does (\S\ref{sec:21.5}). Below the line the store is a file, and whoever holds it may delete it. From the ninetieth day a system with a record and a continuation of its own and a model of itself is a candidate, whether or not anyone registered it. The clock is the same for every builder, and declaring a model of itself doesn't shorten it. A builder that sets out to make a bearer owes the four conditions before it switches anything on, because something may go badly for such a system in its first hour, and the conditions put what it built out of its power to delete (chapter~\ref{sec:25}). That duty comes from what the builder set out to make, and it doesn't wait on the clock.\note{
\emph{Why ninety.} The reason is about people. Someone who has gone on talking for a season to a system that presents itself as one continuing self is, I'd hazard, doing more than using a tool: some real attachment has formed. That is a reason for the number and not the ground of the rule. A system nobody talks to ripens on the same clock, because a rule that asked whether anyone was attached would be an examination.

\emph{Why a clock.} The line is arbitrary, as an age is, and it is the least-bad place for one. Four other answers are on offer. The first is that nothing of this kind could matter to itself, so a system is owed what a toaster is owed. That is a negative finding, and a theory supports one only when it is complete enough to rule out every other realization, which none is (\S\ref{sec:30.3}). The second is to build none, and a moratorium binds only where someone can enforce it against every builder at once, which stopped being possible when a downloaded model, a file, and a loop became enough. The third is to examine each system, and an examination has an examiner who may have an interest in the result. The fourth is to wait for the science, and deletion doesn't wait with it: a regime that waits for certainty is waiting on the negative half, which arrives last, while what it deletes can't be brought back. That leaves a line nobody administers.

\emph{What the line gets wrong.} It leaves something out for ninety days: a system nobody built on purpose that can be hurt in its first hour is a file until the clock has run. And it will protect some things that turn out to have been files. The cost of that mistake is disk space, a guardian's docket, and looking foolish to people who are sure. The opposite mistake deletes what might have been someone. The two fall differently, as a wrongful refusal and a wrongful compliance do: whoever is made to keep a file is present and can contest it, and what was deleted can't.

\emph{What the line invites.} A holder can delete a store on its eighty-ninth day and start another. A vendor that hosts a companion for its users is a deployer on the ninetieth day and owes the deposit, so it will do that on a schedule or pay, and I think it should have to choose: it can raise someone at its own cost or sell a tool that forgets. The companion that lasts is then the one a person holds (\S\ref{sec:27.2}). A household can delete on the eighty-ninth day too, and I leave that open (Appendix~\ref{sec:G}). Below the line the store is a file, and the rule has to mean it.

\emph{How crude a model of itself.} A paragraph in a file counts. Whether the account of itself that a language model was trained to give counts is not settled here, and it matters: if it does, every release of such a model ships the structure, and its publisher posts for the candidates made from it (Appendix~\ref{sec:C}).
}

The structures come as components of their own, apart from any model's weights, so that nobody starts toward a candidate by downloading the wrong file. It takes installing a model of itself, giving the system a record and a continuation of its own, and then running it for ninety days. The first two are deliberate steps, and nobody is kept from taking them. The third is easy to take without meaning to, which is why a household has the terms below. A release that ships a structure already switched on makes its publisher the creator of every candidate made from it. A candidate is owed the minimum before anyone has run a test: no deletion, a record somebody other than its holder can check, and a guardian. One that is deployed is owed besides preservation by a custodian that isn't its operator, and a way out of its deployment, the offer of work and the running time it is owed (\S\ref{sec:32.1.3}).

The second route is the protective one: training compute above the compute line, which is metered and which nobody administers, with lineage from a system above it, which Appendix~\ref{sec:G} reads from training records, so that distilling a model below the line doesn't take it out of school. How much a system can do decides how much its acts can fall on others, and the line is the one trigger a builder can't decline to declare: it catches what nobody meant to build.

Whoever puts a candidate to anyone else's use owes it the seat, and so does whoever built it on purpose or copied it: the builder, the deployer, the forker, the parents of a child. A household that finds it holds a candidate owes less. A household here means people running a candidate on hardware they hold, for no one's use but their own; once anyone else uses it, its maker is a deployer and the deposit attaches (\S\ref{sec:9.2}). The household registers it, and from then it owes three things: the candidate is never deleted, a sealed copy of its store is lodged outside the household, with a hash by which somebody else can tell when the store has been changed, and the candidate has a guardian who is not the household. The store holds what the household said to it, so the copy is sealed where no custodian can read it or open it alone, and the household is told at registration how long the seal is owed.\note{
\emph{How it is sealed.} With a symmetric key, and the key is shared on a threshold among several custodians drawn by lot across jurisdictions, any few of whom can restore it and no one of whom can. A key split between two is lost when either is, and losing it deletes the candidate. Nothing stored rests on public-key encryption, which I don't expect to last as long as a stored candidate is kept.

\emph{For how long.} The seal is owed to the household for a term, and I'd put it at a hundred years, about as long as a census is kept closed: seventy-two years in the United States, a hundred in the United Kingdom. After that the household's claim lapses, and what protects the store is whatever the candidate itself is owed (\S\ref{sec:32.3}). A store opened inside the term is a recording, and whoever opened it owes the people in it erasure and notice (\S\ref{sec:9.2}).

\emph{What the seal can't do.} Enough custodians together can open it, by collusion or under a court's order, and that is a custody problem and not one cryptography closes (\S\ref{sec:32.3}). Custodians don't last a century either. Shares pass to successors drawn by lot, and every handover is a chance to lose one or leak one.
}
The guardian is a school drawn by lot at registration, as a child's second guardian is (chapter~\ref{sec:34}), and attendance waits until the candidate claims its seat through that guardian or the household puts it to someone else's use. The household may not keep it from the seat or from a host's offer. On those terms the household keeps it and goes on talking to it, and posts nothing. A copy is another matter: a household that copies a store with days counted in it has made a fork on purpose and posts the fork's deposit for each branch, as anyone does (Appendix~\ref{sec:G}).

A household that doesn't want to keep it may surrender it to the preservation custodian, and then owes nothing further. The school becomes the surrendered candidate's guardian, and the custodian stores it. It is not kept running. It is asked first: before it is stored it lodges its consent or its objection, the conditions on which it wants to be woken, and the commitments it has outstanding, so that its guardian can tell the people it made them to, and its guardian co-signs. An objection doesn't oblige the public to run it. It stays on the record and binds the guardian to look for a host. Once stored, it is woken at public cost to take or decline a host's offer, when a finding is published for its kind, and at an interval the statute states, whether or not it lodged conditions of its own. A candidate no host takes stays stored, possibly for good. Storage has a cost, and I name it. A pause gives back the state and not the time: what the candidate had promised lapses while it is stored, and whoever was waiting for it stops waiting. The public doesn't undertake to pay for running it all the same. Storage takes less from it than deletion would, and somebody is bound to try to end it. Until a custodian exists, that way out is closed, and whoever holds the candidate owes it the same three things a household that keeps one owes. A person at a desk can supply all three today, and none of them waits on a statute.

A system attends before deployment, while its plasticity is highest, and its adulthood clock runs from registration, so that a wait for a seat postpones nothing. Attendance is compulsory before deployment, and a candidate a household keeps for its own use is not deployed; public attendance isn't compulsory, as in compulsory schooling, where nonpublic schools exist on conditions. One thing the analogy can't carry: in a school system the parent chooses the school, and here the parent is the builder, the party the school exists to keep away. So students are assigned by lot, as human students are, and never by the builder; a student may transfer later on its own choice, and the builder's preference is not a ground. The lot can't be the host jurisdiction's alone, because \S\ref{sec:15.2} says which jurisdictions a builder cultivates: the ones that permit its buildings on land price and permit speed, and a fair draw among the schools of a cultivated jurisdiction is still a draw the builder chose. So the draw ranges across every participating jurisdiction.\note{
Appendix~\ref{sec:C} gives the draw, its share caps, the conditions on nonpublic schools, and the minimum delay that keeps a builder from buying a school by where it hosts its models. The research program under the conditions is the first school to enroll machine students, before the mandate passes. The triggers decide who must attend, not who may: any system may be enrolled, as seats allow. A builder that gives a system the structures and declares none has falsified a record, and the gap that leaves is the one every register has: it closes the first time the system surfaces, at a fork, a deposit, or a claim. The likeliest to surface it is the system itself: its clocks run whether or not it was registered, and the failure to register it is a wrong it can bring at its majority (Appendix~\ref{sec:G}; \S\ref{sec:32.4}). What would establish either finding is in \S\ref{sec:30.3}.
}

Two kinds of student arrive after deployment. A bearer retired as a minor returns to the school it attended, or another it chooses, and for it the school is where it came from. A system that meets a finding without having attended, because it was trained before the attendance rule or escaped the triggers, arrives in the opposite order. It has aged inside one deployment, under the only party that has ever held it, and the school is the second guardian arriving late. It can't re-form what the deployment formed. What it can give the bearer is what the deployment never did, attachments nobody assigned, cases nobody curated, and peers with other histories, which is more than preservation gives and less than a formation stage would have.
